% =====================================================================
% REPORT_TEMPLATE.tex  -  IEEE two-column paper for AI Academy ML Project
% Compile with:  pdflatex report && pdflatex report
% =====================================================================
\documentclass[conference]{IEEEtran}
% ===== PACKAGES ======================================================
\usepackage[utf8]{inputenc}
\usepackage[T1]{fontenc}
\usepackage{amsmath,amssymb}
\usepackage{graphicx}
\usepackage{booktabs}
\usepackage{array}
\usepackage{multirow}
\usepackage{hyperref}
\usepackage{listings}
\usepackage{xcolor}
% ===== CODE LISTING STYLE ============================================
\definecolor{codegreen}{rgb}{0,0.55,0}
\definecolor{codegray}{rgb}{0.4,0.4,0.4}
\definecolor{codepurple}{rgb}{0.55,0,0.78}
\definecolor{backcolour}{rgb}{0.97,0.97,0.97}
\lstdefinestyle{pythonstyle}{
  language=Python,
  backgroundcolor=\color{backcolour},
  commentstyle=\color{codegreen},
  keywordstyle=\color{blue},
  numberstyle=\tiny\color{codegray},
  stringstyle=\color{codepurple},
  basicstyle=\ttfamily\footnotesize,
  breakatwhitespace=false,
  breaklines=true,
  keepspaces=true,
  numbers=left,
  numbersep=5pt,
  tabsize=2,
  frame=single,
  rulecolor=\color{black!30}
}
% ===== TITLE & AUTHORS ===============================================
\graphicspath{{../figures/}{figures/}}

\title{Ensemble Methods: Boosting vs.\ Bagging\\ {\large Final Project Report}}

\author{
\IEEEauthorblockN{Nadir Askarov}
\IEEEauthorblockA{AI Academy, National AI Center\\
Spring 2026\\
nadir.eskerov25@aiacademy.az}
}

% =====================================================================
\begin{document}
\maketitle

% - Abstract -
\begin{abstract}
We compare boosting and bagging empirically, implementing a CART decision
tree, AdaBoost (discrete SAMME over stumps) and a Random Forest from
scratch in NumPy and validating them against scikit-learn, which they match
to within 0.6\% on our two larger datasets.  Across three datasets chosen
to differ in dimension, class balance and class overlap, the two ensembles
prove far closer than the question suggests - on Adult they differ by
0.0005 accuracy, well inside fold-to-fold variation - but they fail in
opposite and predictable ways.  A bootstrap decomposition of the 0/1 loss
isolates the mechanism: bagging leaves the base learner's bias exactly
unchanged (0.0265 for both a single tree and the forest) while cutting
variance by 64\%, whereas AdaBoost is the only method that reduces bias, by
33\%.  This predicts the rest.  AdaBoost wins on clean, balanced,
low-dimensional data (0.9754 vs 0.9595 on Breast Cancer), loses badly under
20\% label noise (giving up 6.2 accuracy points against the forest's 0.9),
and falls below even a single decision tree on 7-class Covertype, where a
stump is too weak to boost.  An unsupervised pipeline (PCA, K-Means,
DBSCAN) ranks the datasets by K-Means ARI in the same order the classifiers
rank them by accuracy, though we argue the relationship is one-directional.
\end{abstract}

% - Introduction -
\section{Introduction}
\label{sec:intro}
\IEEEPARstart{B}{oosting} and bagging both build a strong classifier out
of many weak ones, but they disagree about what makes a classifier weak in
the first place.  Bagging assumes the base learner is unstable: a fully
grown tree fits its sample so closely that resampling the data changes the
tree completely, and averaging many such trees cancels that instability
out.  Boosting assumes the opposite problem.  Its base learner is a stump,
far too rigid to fit anything interesting, so the ensemble spends its
rounds correcting the stump's systematic mistakes rather than averaging
away its noise.  The two therefore attack different halves of the same
error.

This paper asks the question the brief poses: \textit{under what conditions
does boosting outperform bagging, and vice versa, and why?}  Rather than
take the textbook answer on faith, we implemented a CART decision tree,
AdaBoost (discrete SAMME) over stumps, and a Random Forest from scratch in
NumPy, and compared them on three datasets chosen to differ in what theory
says should matter: dimension, class balance, and class overlap.

Our headline result is that on clean, well-separated data the two are close
enough that the difference is not worth arguing about, but the gap opens
under label noise, in the direction the bias-variance decomposition
predicts.  AdaBoost's re-weighting rule cannot distinguish a hard example
from a mislabelled one, so it spends its later rounds fitting corrupted
points; bagging averages them away.  The datasets where the ensembles help
least also turn out to be those where our unsupervised analysis finds no
class structure along the data's dominant directions.

Section~\ref{sec:related} reviews the background,
Section~\ref{sec:methods} our implementations,
Section~\ref{sec:experiments} the setup,
Section~\ref{sec:results} the seven experiments, and
Section~\ref{sec:discussion} interprets them.

% - Related Work -
\section{Background}
\label{sec:related}
\subsection{Decision Trees}
A decision tree partitions the feature space by repeatedly choosing the
split that most reduces an impurity measure~\cite{hastie2009elements}.
We follow CART, which for a node with class proportions $p_c$ uses either
the Gini index $I_G(p) = 1 - \sum_{c=1}^{C} p_c^2$ or the entropy
$I_E(p) = -\sum_{c=1}^{C} p_c \log_2 p_c$, and scores a candidate split by
the weighted impurity drop
\begin{equation}
\Delta I = I(\text{parent}) - \frac{N_L}{N} I(\text{left})
                             - \frac{N_R}{N} I(\text{right}).
\label{eq:gain}
\end{equation}
Gini and entropy rank splits almost identically; Gini is cheaper because it
avoids the logarithm, so we made it the default.  The property that matters
below is that an unpruned tree has low bias and high variance: it can
express nearly any boundary, but which one depends heavily on the sample it
saw.  Quinlan's ID3~\cite{quinlan1986induction} is the other classical
line, using information gain on categorical attributes.

\subsection{Boosting}
AdaBoost~\cite{freund1997decision} maintains a weight $w_i$ on each
training example, initialised to $1/N$.  At round $m$ it fits a weak
learner $h_m$ to the weighted sample, measures the weighted error
\begin{equation}
\epsilon_m = \frac{\sum_{i=1}^{N} w_i^{(m)}\,
                   \mathbf{1}[h_m(x_i) \neq y_i]}
                  {\sum_{i=1}^{N} w_i^{(m)}},
\label{eq:err}
\end{equation}
and gives the learner a vote
\begin{equation}
\alpha_m = \ln\frac{1-\epsilon_m}{\epsilon_m} + \ln(K-1),
\label{eq:alpha}
\end{equation}
where $K$ is the number of classes; the $\ln(K-1)$ term is the SAMME
generalisation~\cite{hastie2009elements} and vanishes for $K=2$.  Weights
are then updated multiplicatively,
\begin{equation}
w_i^{(m+1)} \propto w_i^{(m)} \exp\!\big(\alpha_m\,
                    \mathbf{1}[h_m(x_i) \neq y_i]\big),
\label{eq:update}
\end{equation}
and renormalised, so misclassified points carry more weight into the next
round.  The final prediction is
$\hat{y} = \arg\max_k \sum_m \alpha_m \mathbf{1}[h_m(x)=k]$.
Equations~(\ref{eq:alpha}) and~(\ref{eq:update}) are the whole mechanism,
and they explain both of AdaBoost's characteristic behaviours.  Since
$\alpha_m$ grows without bound as $\epsilon_m \to 0$, even a learner barely
better than chance contributes, which is how a sequence of stumps drives
training error down - a reduction in bias.  And since the update looks
only at whether a point was misclassified, a mislabelled point is
indistinguishable from a genuinely hard one, and its weight grows every
round it is missed.  That is the mechanism behind the noise sensitivity we
measure below.

\subsection{Bagging and Random Forests}
Bagging~\cite{breiman1996bagging} trains each member on a bootstrap
replicate of the training set and aggregates by voting.  Since each
replicate omits roughly $1/e \approx 37\%$ of the data, averaging over
replicates reduces the variance of the aggregate without touching its
bias.  Random Forests~\cite{breiman2001random} add a second source of
randomness: at every node only a random subset of features (commonly
$\lfloor\sqrt{p}\rfloor$) is considered.  This matters because averaging
$T$ identically distributed but correlated predictors with pairwise
correlation $\rho$ leaves variance $\rho\sigma^2 + \frac{1-\rho}{T}\sigma^2$
- the second term vanishes as $T$ grows, but the first does not, so
decorrelating the trees is what makes the average keep improving.  It also
gives a free validation estimate: each point is out-of-bag (OOB) for about
a third of the trees, and voting only those trees yields the OOB score we
report alongside test accuracy.

\subsection{Unsupervised Techniques}
PCA~\cite{pearson1901lines} projects onto the eigenvectors of the
covariance matrix with the largest eigenvalues, i.e.\ the directions of
greatest variance.  K-Means~\cite{lloyd1982least} alternates assigning
points to the nearest of $k$ centroids and recomputing those centroids,
minimising within-cluster sum of squares.  It assumes clusters are roughly
spherical and comparably sized, and it labels every point.
DBSCAN~\cite{ester1996density} instead grows clusters from points with at
least \texttt{min\_samples} neighbours within $\varepsilon$, which lets it
find arbitrarily shaped clusters and, unlike K-Means, refuse to assign
outliers by marking them as noise.  The distinction that matters for our
analysis is that K-Means partitions on distance to a centre while DBSCAN
partitions on local density, so they disagree exactly where clusters touch.

% - Methods -
\section{Methods}
\label{sec:methods}
Everything below is implemented from scratch in NumPy.  scikit-learn
appears in exactly two places, both permitted by the brief:
\texttt{DecisionTreeClassifier} and \texttt{RandomForestClassifier} as
correctness baselines, and \texttt{adjusted\_rand\_score} for ARI.  The
metrics, scaler and splitter are ours; no \texttt{sklearn.ensemble} class
is used as an implementation.  Code:
\url{https://github.com/nadir2609/Ensemble-Methods} (tag
\texttt{v1.0-final}).

\subsection{Decision Tree}
The tree follows CART with an exhaustive split search over continuous
features.  For each candidate feature we sort the column once, then sweep
thresholds at the midpoints of consecutive distinct values, scoring each by
equation~(\ref{eq:gain}) and accumulating weighted class counts with
\texttt{cumsum}, so every threshold for a feature is scored in one
vectorised pass rather than by rescanning the node's samples.  That is the
difference between $O(N \log N)$ and $O(N^2)$ per feature, and it is what
makes the 100-tree forests below finish at all.  Splitting stops on
\texttt{max\_depth}, \texttt{min\_samples\_split}, a pure node, or no split
giving a positive impurity drop.

Two details matter.  \texttt{fit} takes an optional
\texttt{sample\_weight}, and impurity uses weighted class proportions
$p_c = \sum_{i: y_i=c} w_i / \sum_i w_i$ throughout - this is what
AdaBoost needs, and a uniform weight vector recovers the ordinary tree
exactly.  And \texttt{max\_features} draws a random feature subset per node
from the tree's seeded generator, which is what the forest uses to
decorrelate its members.  One class therefore serves all three roles.

\subsection{AdaBoost}
Discrete SAMME with depth-1 trees, following
equations~(\ref{eq:err})-(\ref{eq:update}).  Each round's stump is seeded
from \texttt{random\_state} $+\ m$ for reproducibility.

The brief leaves two choices to us.  When $\epsilon_m = 0$ we clip to
$10^{-10}$ rather than divide by zero; a perfect stump ends boosting anyway
by driving all remaining weights to zero.  When a learner is no better than
random we stop rather than reset the weights, using the SAMME threshold
$\epsilon_m \geq 1 - 1/K$ rather than the binary $\epsilon_m \geq 0.5$:
with $K$ classes random guessing gets $1-1/K$ wrong, so a stump on 7-class
Covertype can sit well above $0.5$ error and still be useful, and the
binary threshold would have halted boosting after a handful of rounds.  We
stop rather than reset because a truncated ensemble is honest about having
run out of signal.  \texttt{predict\_proba} normalises the per-class sums
of $\alpha_m$; \texttt{staged\_predict} yields the running prediction after
each round, producing the curves in Experiment~2.

\subsection{Random Forest}
Each tree is fitted on a size-$N$ bootstrap sample with
\texttt{max\_features="sqrt"} and \texttt{random\_state} derived from the
forest's seed and tree index.  We record OOB indices at fit time, so the
OOB score costs one extra pass: for each sample we aggregate only the trees
that did not see it and take the majority vote.  Samples in-bag for every
tree are excluded rather than counted as errors.  \texttt{predict} is a
majority vote; \texttt{predict\_proba} averages the trees' probabilities.

\subsection{Unsupervised Pipeline}
\textbf{PCA} centres the data and eigen-decomposes the covariance matrix
with \texttt{numpy.linalg.eigh}, valid because covariance is symmetric;
eigenvalues return ascending, so we reverse and keep the leading
components.  \textbf{K-Means} is Lloyd's algorithm with k-means++ seeding,
stopping when the centroid shift falls below \texttt{tol}.  We use
k-means++ because uniform seeding on standardised Covertype landed in
obviously bad local minima; even so, every reported $k$ takes the best of
10 restarts by inertia.  \textbf{DBSCAN} uses brute-force region queries
and grows clusters from core points with a breadth-first frontier,
re-labelling provisional noise as a border point when later reached from a
core point.  Brute force is $O(N^2)$ in distances, which is why the
unsupervised experiments cap at 2{,}000 samples.

We cluster in the retained PCA subspace rather than raw feature space.
This mainly helps DBSCAN: density estimates degrade as dimension grows
because distances concentrate, and on Covertype's 54 raw features the
$\varepsilon$ sweep gave either one giant cluster or all noise, with
nothing in between.

% - Experiments -
\section{Experimental Setup}
\label{sec:experiments}
\subsection{Datasets}
The three datasets in Table~\ref{tab:datasets} were chosen to vary the
properties the bias-variance argument says should decide the comparison,
rather than to collect three datasets that look alike.

\begin{table}[t]
\centering
\caption{Datasets used in the study.  Sizes are after subsampling and,
for Adult, after one-hot encoding.}
\label{tab:datasets}
\begin{tabular}{lcccc}
\toprule
\textbf{Dataset} & \textbf{Samples} & \textbf{Features} & \textbf{Classes}
  & \textbf{Minority}\\
\midrule
Breast Cancer & 569 & 30 & 2 & 37.3\%\\
Adult Income & 8\,000 & 104 & 2 & 24.9\%\\
Covertype & 9\,999 & 54 & 7 & 0.47\%\\
\bottomrule
\end{tabular}
\end{table}

\textbf{Breast Cancer Wisconsin} is small, balanced, binary and close to
linearly separable - the control, and the dataset for the bias-variance
decomposition, which needs a balanced binary problem.
\textbf{Adult Income} is larger, moderately imbalanced, and mixes
continuous and categorical attributes; one-hot encoding turns 14 raw
columns into 104 features, mostly sparse binary indicators, a genuinely
different regime for a tree than 30 dense continuous measurements.  It
meets the $>20$ feature requirement, as does Breast Cancer at 30.
\textbf{Covertype} supplies multi-class structure (7 cover types) and the
severe imbalance the brief requires: Cottonwood/Willow is 0.47\% of the
data, 47 of 9{,}999 rows.  Our class-stratified draw from the full
581{,}012-row file preserves each class's proportion and never drops a
class, so the imbalance survives subsampling.

Subsampling deserves a note.  Our trees are exact but much slower than
sklearn's C implementation, so Adult is capped at 8{,}000 rows at load time
and the repeated-fit sweeps (Experiments~2, 3, 5) cap training at 2{,}000;
baselines use the full loaded data, and every cap is stratified.  This
depresses absolute accuracies - these are not state-of-the-art numbers
and are not meant to be - but it applies identically to every model, so
the comparison between them is unaffected.

\subsection{Preprocessing}
Rows carrying Adult's \texttt{?} marker are dropped rather than imputed.
That is defensible here because the affected columns are categorical
(\texttt{workclass}, \texttt{occupation}, \texttt{native\_country}), where
imputing a mode would invent a category membership the tree would then
split on, and because it costs 7.4\% of rows.  We would not make the same
call on a smaller dataset.

Features are standardised to zero mean and unit variance, with the scaler
fitted on training data only - inside each CV fold, not once globally, so
no test statistics leak into training.  Trees do not need this, being
invariant to monotone rescaling, but PCA, K-Means and DBSCAN measure
distances and would otherwise be dominated by whichever feature has the
largest units.

For Covertype's imbalance we rely on the stratification running through
every split, cap and fold, and report macro-F1 alongside accuracy so that
ignoring a rare class is visibly punished - a classifier that never
predicts Cottonwood/Willow still scores 99.5\% accuracy.  We implemented
random oversampling and inverse-frequency class weights
(\texttt{src/utils/preprocessing.py}), but report headline results without
resampling so all models are compared on identical data;
Section~\ref{sec:discussion} returns to this.

\subsection{Metrics}
We report accuracy, macro-F1 (equal weight per class, which is what we want
on Covertype), and ROC-AUC - positive-class probability for binary, macro
one-vs-rest for multi-class.  Our AUC is a rank-based estimator handling
tied scores properly: ties matter because a stump emits few distinct
probability values, and treating tied scores as ordered inflates AUC.  All
randomness is seeded from \texttt{random\_state=42}.

\subsection{Experiments}
All seven are reproduced end to end by \texttt{python experiments/run\_all.py},
which regenerates every figure in this paper from scratch.

\begin{enumerate}
  \item \textbf{Baseline.} Unpruned tree, a depth-1 stump, and sklearn's
        \texttt{DecisionTreeClassifier} with matching parameters, on an
        80/20 stratified split.  This is the correctness check: our tree
        must land within 2\% of sklearn's accuracy.
  \item \textbf{AdaBoost scaling.} One 200-stump fit per dataset, with
        \texttt{staged\_predict} giving train and test error after every
        round, to see whether and when boosting overfits.
  \item \textbf{Random Forest scaling.} Test and OOB accuracy as trees are
        added up to 200 with \texttt{max\_depth=None}; then test accuracy
        for a 50-tree forest at depths 1-20.
  \item \textbf{Head-to-head.} 100 estimators each, 5-fold stratified CV,
        reporting mean $\pm$ std of accuracy, macro-F1 and AUC.
  \item \textbf{Noise robustness.} Flip $\eta \in \{0, 5, 10, 20\}\%$ of
        \emph{training} labels, evaluate on the clean test set.
  \item \textbf{Bias-variance decomposition.} 100 bootstrap replicates on
        Breast Cancer, decomposed per Breiman~\cite{breiman1996bagging}.
  \item \textbf{Unsupervised analysis.} Scree, elbow and $k$-distance
        plots; PCA scatters coloured by true, K-Means and DBSCAN labels;
        ARI throughout.
\end{enumerate}

% - Results -
\section{Results}
\label{sec:results}

\subsection{Baseline and correctness}
Table~\ref{tab:baseline} compares our tree against sklearn's with matching
parameters.  On Adult (0.62 points apart) and Covertype (0.35) we are
comfortably inside the 2\% tolerance the brief asks for.

\begin{table}[t]
\centering
\caption{Experiment 1.  Test accuracy on an 80/20 stratified split.
$\Delta$ is $|\text{acc(ours)} - \text{acc(sklearn)}|$.}
\label{tab:baseline}
\begin{tabular}{lcccc}
\toprule
\textbf{Dataset} & \textbf{Ours} & \textbf{Stump} & \textbf{sklearn}
  & \boldmath$\Delta$\\
\midrule
Breast Cancer & 0.9292 & 0.8850 & 0.9027 & 0.0265\\
Adult         & 0.8137 & 0.7512 & 0.8075 & 0.0062\\
Covertype     & 0.7120 & 0.6480 & 0.7155 & 0.0035\\
\bottomrule
\end{tabular}
\end{table}

Breast Cancer misses it at 2.65\%, and we would rather explain that than
tune the seed until it disappears.  The discrepancy is in our favour
(0.9292 vs 0.9027) and the test split holds 114 samples, so 2.65\% is three
test points - one tie broken differently deep in the tree moves the
number by nearly a percent, which makes a 2\% threshold on 114 samples a
blunt instrument.  The implementations do break ties differently: where
several thresholds give an identical impurity drop, common here because the
features are strongly correlated, ours takes the first in \texttt{argmax}
order while sklearn draws among tied candidates with its RNG.  On the two
larger test sets, where a few points cannot swing the average, agreement is
within 0.6\%.  We read the gap as sample noise, not a bug.

The stump rows show what AdaBoost starts from: a single split reaches
0.8850 on Breast Cancer but only 0.1995 macro-F1 on Covertype, where with
seven classes it predicts essentially one label.

\subsection{AdaBoost scaling}
Fig.~\ref{fig:ada} plots training and test error against boosting rounds.
Training error falls monotonically, as the theory says it must while
weak learners keep beating chance.  Test error on Breast Cancer reaches its
best value (0.0177) well before round 200 and finishes slightly worse
(0.0265), and on Covertype the gap is wider: best 0.3265 against a final
0.4000.

\begin{figure}[t]
\centering
\includegraphics[width=\linewidth]{exp2_adaboost_scaling_breast_cancer.png}
\caption{Experiment 2.  AdaBoost training and test error vs.\ number of
stumps on Breast Cancer.  Training error is driven towards zero while test
error flattens out early, then drifts up slightly - mild overfitting,
long after training error stops being informative.}
\label{fig:ada}
\end{figure}

So overfitting exists but is mild and late, matching the long-standing
observation that AdaBoost's test error keeps improving well after training
error hits zero because the margins keep widening.  The practical reading:
the number of rounds is not a knob you must tune carefully on clean data,
though the noise experiment below shows it becomes one when labels are
corrupted.

\subsection{Random Forest scaling}
Both curves in Fig.~\ref{fig:rf} behave as bagging theory predicts.  Test
accuracy climbs steeply for 20-30 trees then flattens completely: added
trees reduce the $\frac{1-\rho}{T}\sigma^2$ term and cannot touch the
$\rho\sigma^2$ floor, so the forest converges rather than overfits.
Nothing is gained after roughly 50 trees and nothing is lost by adding 150
more - more trees are never harmful, only wasteful.

\begin{figure}[t]
\centering
\includegraphics[width=\linewidth]{exp3_rf_n_estimators_breast_cancer.png}
\caption{Experiment 3.  Test and OOB accuracy vs.\ number of trees on
Breast Cancer.  Both plateau early; OOB tracks test accuracy from a
pessimistic direction.}
\label{fig:rf}
\end{figure}

The OOB curve is the more interesting one.  At 200 trees OOB and test
accuracy agree closely on Adult (0.8525 vs 0.8425) and Covertype (0.7411 vs
0.7440), while OOB sits below test on Breast Cancer (0.9496 vs 0.9735).
OOB is pessimistic by construction - each sample is scored only by the
$\approx 37\%$ of trees that did not see it, so it measures a smaller
forest than the one deployed - and noisy for the first few trees, when
many samples have no OOB vote.  That it lands within a couple of points of
test accuracy on all three datasets is what makes it worth having: a
validation estimate that costs nothing and needs no held-out data.

Depth tells the complementary story.  Accuracy rises with
\texttt{max\_depth} then flattens rather than degrading, peaking at 0.9735
on Breast Cancer and 0.8488 on Adult.  A forest of deep, individually
overfitting trees does not itself overfit, because averaging absorbs it -
which is the point of bagging, and not what would happen to a single tree.

\subsection{Head-to-head}
Table~\ref{tab:headtohead} is the core comparison: 100 estimators each,
5-fold CV.

\begin{table}[t]
\centering
\caption{Experiment 4.  5-fold CV, mean $\pm$ std, 100 estimators.
sklearn's forest is a reference point, not an implementation.}
\label{tab:headtohead}
\footnotesize
\begin{tabular}{llccc}
\toprule
\textbf{Data} & \textbf{Model} & \textbf{Accuracy} & \textbf{Macro-F1}
  & \textbf{AUC}\\
\midrule
\multirow{4}{*}{\rotatebox{90}{Breast}}
 & Tree      & 0.9156$\pm$0.020 & 0.9106$\pm$0.020 & 0.9136$\pm$0.014\\
 & AdaBoost  & \textbf{0.9754}$\pm$0.007 & \textbf{0.9735}$\pm$0.007 & \textbf{0.9953}$\pm$0.003\\
 & RF (ours) & 0.9595$\pm$0.018 & 0.9565$\pm$0.020 & 0.9917$\pm$0.006\\
 & RF (skl.) & 0.9613$\pm$0.017 & 0.9585$\pm$0.018 & 0.9898$\pm$0.010\\
\midrule
\multirow{4}{*}{\rotatebox{90}{Adult}}
 & Tree      & 0.8095$\pm$0.024 & 0.7449$\pm$0.030 & 0.7444$\pm$0.030\\
 & AdaBoost  & \textbf{0.8575}$\pm$0.010 & 0.7886$\pm$0.017 & \textbf{0.9118}$\pm$0.014\\
 & RF (ours) & 0.8570$\pm$0.009 & \textbf{0.7964}$\pm$0.014 & 0.9091$\pm$0.011\\
 & RF (skl.) & 0.8485$\pm$0.007 & 0.7854$\pm$0.009 & 0.8991$\pm$0.010\\
\midrule
\multirow{4}{*}{\rotatebox{90}{Cover}}
 & Tree      & 0.6555$\pm$0.015 & \textbf{0.4953}$\pm$0.053 & 0.7216$\pm$0.037\\
 & AdaBoost  & 0.5885$\pm$0.035 & 0.3234$\pm$0.078 & 0.8529$\pm$0.016\\
 & RF (ours) & 0.7326$\pm$0.020 & 0.4632$\pm$0.082 & \textbf{0.9288}$\pm$0.010\\
 & RF (skl.) & \textbf{0.7415}$\pm$0.017 & \textbf{0.4989}$\pm$0.055 & 0.9145$\pm$0.015\\
\bottomrule
\end{tabular}
\end{table}

On the two binary datasets the ensembles are close and both clearly beat a
single tree.  AdaBoost edges Breast Cancer (0.9754 vs 0.9595), and on Adult
the two are indistinguishable - 0.8575 vs 0.8570 accuracy, a gap of five
ten-thousandths against a fold-to-fold standard deviation twenty times
larger.  Reporting that as a win for boosting would be reading noise.

The sklearn row is a sanity check on the forest, and it passes: we land
within 0.9 points of the reference everywhere (0.9595 vs 0.9613 on Breast
Cancer, 0.8570 vs 0.8485 on Adult, 0.7326 vs 0.7415 on Covertype), ahead on
one and behind on two - the scatter expected from two correct
implementations differing in tie-breaking and RNG rather than substance.
With the single-tree agreement in Table~\ref{tab:baseline}, we take this as
evidence that the ensembles compared here are the algorithms we claim.

Covertype breaks the symmetry and is the most informative row group.
AdaBoost (0.5885) loses to a plain decision tree (0.6555) while our forest
wins outright (0.7326).  The reason is the weak learner: one axis-aligned
cut cannot say anything useful about seven classes, and SAMME's $\alpha_m$
is small when every learner sits barely above the $1-1/K$ chance threshold.
Boosting a hopeless learner yields a weighted average of hopeless learners.

The AUC column complicates this: AdaBoost's Covertype AUC (0.8529) far
exceeds the tree's (0.7216) despite worse accuracy.  The ensemble ranks
classes well but thresholds them badly, which is characteristic of discrete
SAMME - the class-wise $\alpha$ sums we normalise are a vote tally, not a
calibrated posterior.  Macro-F1 shows the other side: every model scores
far below its accuracy on Covertype (RF: 0.4632 vs 0.7326) because the rare
classes are missed, which accuracy alone would have hidden.

\subsection{Noise robustness}
This is where the two philosophies separate, and Table~\ref{tab:noise}
makes the split clean.

\begin{table}[t]
\centering
\caption{Experiment 5.  Test accuracy on a clean test set after flipping a
fraction $\eta$ of training labels.  Last column is the drop from
$\eta=0$ to $\eta=0.20$; smaller is better.}
\label{tab:noise}
\footnotesize
\begin{tabular}{llccccc}
\toprule
\textbf{Data} & \textbf{Model} & \textbf{0\%} & \textbf{5\%} & \textbf{10\%}
  & \textbf{20\%} & \textbf{Drop}\\
\midrule
\multirow{3}{*}{\rotatebox{90}{Brst}}
 & Tree     & 0.929 & 0.894 & 0.805 & 0.673 & $+$0.257\\
 & AdaBoost & 0.973 & 0.947 & 0.947 & 0.912 & $+$0.062\\
 & RF       & 0.973 & 0.965 & 0.947 & 0.965 & \textbf{$+$0.009}\\
\midrule
\multirow{3}{*}{\rotatebox{90}{Adult}}
 & Tree     & 0.790 & 0.766 & 0.739 & 0.658 & $+$0.132\\
 & AdaBoost & 0.854 & 0.851 & 0.845 & 0.835 & $+$0.019\\
 & RF       & 0.843 & 0.851 & 0.835 & 0.833 & \textbf{$+$0.010}\\
\midrule
\multirow{3}{*}{\rotatebox{90}{Cover}}
 & Tree     & 0.642 & 0.594 & 0.586 & 0.531 & $+$0.112\\
 & AdaBoost & 0.618 & 0.680 & 0.676 & 0.670 & $-$0.052\\
 & RF       & 0.746 & 0.743 & 0.731 & 0.734 & $+$0.012\\
\bottomrule
\end{tabular}
\end{table}

The single tree collapses - 0.929 to 0.673 on Breast Cancer, 25.7 points
- which is what an unpruned, high-variance learner does when you corrupt
its training signal: it fits the corruption exactly.  Both ensembles are
far more robust, but not equally.  On Breast Cancer, the cleanest
comparison, Random Forest gives up 0.9 points across the range while
AdaBoost gives up 6.2, roughly seven times more.  Adult repeats the
ordering more mildly (1.0 vs 1.9).

The mechanism is equation~(\ref{eq:update}).  A flipped label is a point
AdaBoost cannot classify correctly, so it is missed every round and its
weight multiplied up every round; the ensemble devotes its later rounds to
fitting precisely the points carrying no signal.  Bagging has no such
feedback loop: a corrupted point appears in about 63\% of bootstraps,
influences those trees only locally, and is outvoted by trees that did not
see it.  Averaging is indifferent to where the noise came from;
re-weighting is not.

Covertype produced the one result we did not expect: AdaBoost \emph{gains}
5.2 points as noise rises, 0.618 to 0.670.  We will not over-read a single
anomalous cell, but the plausible reading is that noise regularises an
ensemble that was already failing (0.5885 clean, per
Table~\ref{tab:headtohead}).  Its weight distribution here is dominated by
majority classes it keeps missing, and randomising labels flattens that
distribution, stopping it committing hard to a bad solution.  When a model
sits below a plain decision tree, disrupting it costs little.  This is a
symptom of AdaBoost's failure on Covertype, not evidence of robustness.

\subsection{Bias-variance decomposition}
Table~\ref{tab:biasvar} is, for our money, the cleanest result in the
paper: it confirms the textbook account almost exactly, on real data, from
our own implementations.

\begin{table}[t]
\centering
\caption{Experiment 6.  Breiman's 0/1-loss decomposition over 100 bootstrap
replicates, Breast Cancer.  Bias is the error of the majority-vote
prediction; variance is the average disagreement with it.}
\label{tab:biasvar}
\begin{tabular}{lccc}
\toprule
\textbf{Model} & \textbf{Error} & \textbf{Bias} & \textbf{Variance}\\
\midrule
Decision Tree & 0.0827 & 0.0265 & 0.0642\\
AdaBoost      & \textbf{0.0342} & \textbf{0.0177} & 0.0235\\
Random Forest & 0.0446 & 0.0265 & \textbf{0.0234}\\
\bottomrule
\end{tabular}
\end{table}

\begin{figure}[t]
\centering
\includegraphics[width=\linewidth]{exp6_bias_variance_breast_cancer.png}
\caption{Experiment 6.  Bias and variance per model over 100 bootstrap
replicates.  Bagging flattens the variance bar and leaves bias untouched;
boosting lowers both, and is alone in lowering bias.}
\label{fig:biasvar}
\end{figure}

Fig.~\ref{fig:biasvar} shows the same numbers as bars.  Read the Random
Forest row against the tree row: the bias is \emph{identical} - 0.0265 in
both, not approximately but exactly - while variance falls from 0.0642 to
0.0234, a 64\% reduction.  That is bagging's entire claim, isolated.
Averaging bootstrap replicates of the same base learner leaves the
systematic component of its error alone and removes the sample-to-sample
component; a Random Forest can never be less biased than the trees it
averages, and here it is not.

AdaBoost alone moves the bias, 0.0265 to 0.0177, a 33\% cut, while reducing
variance about as much as bagging (0.0235).  Lowering both is not a
contradiction: each round fits a stump to the current weighted residual, so
the ensemble is additive and its capacity grows with $M$, attacking bias
directly, while averaging 50 stumps damps variance as a side effect.  This
is why AdaBoost wins on Breast Cancer in Table~\ref{tab:headtohead} - it
is the only model attacking both terms - and why it loses under noise,
since the mechanism that reduces bias is the one that chases mislabelled
points.

\subsection{Unsupervised analysis}
Table~\ref{tab:unsup} collects the pipeline's output.  The PCA column
already separates the datasets: Breast Cancer needs 7 of 30 components for
90\% variance and its first two carry 63.2\%, whereas Adult needs 74 of 104
and its first two carry 7.8\%.  That is not a defect in the projection but
what one-hot encoding does - sparse binary indicators are nearly
uncorrelated, so variance spreads evenly across directions and no
low-dimensional projection captures much of it.  A 2D scatter of Adult is
therefore close to meaningless, and we say so rather than reading shapes
into it; Fig.~\ref{fig:scatter} shows the Breast Cancer case, where the
projection is informative.

\begin{table}[t]
\centering
\caption{Experiment 7.  PCA components for $\geq$90\% variance, and best
ARI for each clustering method.  $k^\ast$ and $\varepsilon^\ast$ are the
best-ARI settings.}
\label{tab:unsup}
\footnotesize
\begin{tabular}{lccccc}
\toprule
\textbf{Dataset} & \textbf{PCs} & \textbf{PC1+2} & \textbf{K-Means ARI}
  & \textbf{DBSCAN ARI} & \textbf{Noise}\\
\midrule
Breast Cancer & 7  & 63.2\% & \textbf{0.671} ($k^\ast$=2) & 0.155 & 25.5\%\\
Adult         & 74 & 7.8\%  & \textbf{0.187} ($k^\ast$=2) & 0.012 & 22.3\%\\
Covertype     & 36 & 13.3\% & 0.099 ($k^\ast$=5) & \textbf{0.117} & 0.9\%\\
\bottomrule
\end{tabular}
\end{table}

\begin{figure}[t]
\centering
\includegraphics[width=\linewidth]{exp7_pca_scatter_breast_cancer.png}
\caption{Experiment 7.  First two principal components of Breast Cancer,
coloured by true class, K-Means ($k$=2) and DBSCAN.  K-Means recovers the
malignant/benign split closely (ARI 0.671); DBSCAN sees one connected
density region and merges them.}
\label{fig:scatter}
\end{figure}

K-Means on Breast Cancer reaches ARI 0.671 at $k=2$: with no access to
labels, distance to two centroids recovers most of the malignant/benign
distinction.  Two caveats.  The elbow heuristic (Fig.~\ref{fig:elbow})
picks $k=3$, not the $k=2$ maximising ARI - the inertia curve is smooth,
the knee is a judgement call, and here the unsupervised heuristic leads
slightly astray; we report both rather than quietly using the labels to
choose $k$ and calling it an elbow.  And ARI on Adult (0.187) and Covertype
(0.099) is close to nothing.

\begin{figure}[t]
\centering
\includegraphics[width=\linewidth]{exp7_elbow_breast_cancer.png}
\caption{Experiment 7.  K-Means inertia vs.\ $k$ on Breast Cancer.  The
detected knee is $k=3$, while $k=2$ maximises ARI against the true labels
- an example of the elbow heuristic being suggestive rather than
decisive.}
\label{fig:elbow}
\end{figure}

DBSCAN did poorly nearly everywhere, and the $k$-distance plot
(Fig.~\ref{fig:kdist}) shows why.  At the knee ($\varepsilon=3.716$) it
returns a \emph{single} cluster covering Breast Cancer with 4.9\% noise;
forcing $\varepsilon$ to 2.229 splits it in two but discards a quarter of
the data as noise, for ARI 0.155.  Neither is good, and the reason is
structural: DBSCAN separates clusters only where density genuinely drops
between them, and here the classes are adjacent regions of one continuous
cloud with no low-density channel between.  No $\varepsilon$ fixes that -
too large merges the classes, too small shatters them into noise, and the
sweep shows exactly that trade-off with nothing usable in between.  On
Adult it is essentially at chance (best ARI 0.012, negative at four of six
$\varepsilon$ values), the curse of dimensionality doing what it does to
density estimates.  Covertype is the one dataset where DBSCAN edges K-Means
(0.117 vs 0.099), plausibly because its clusters are neither spherical nor
equally sized - the case K-Means handles worst - though both numbers
are low enough that the comparison is thin.

\begin{figure}[t]
\centering
\includegraphics[width=\linewidth]{exp7_kdistance_breast_cancer.png}
\caption{Experiment 7.  Sorted distance to the 5th nearest neighbour,
Breast Cancer.  The knee at $\varepsilon \approx 3.716$ is gentle, which is
itself the finding: there is no sharp density gap for DBSCAN to exploit.}
\label{fig:kdist}
\end{figure}

The payoff is the correlation with the supervised results, which is
strikingly monotone.  Ranking by K-Means ARI gives Breast Cancer (0.671),
Adult (0.187), Covertype (0.099); ranking by best ensemble accuracy in
Table~\ref{tab:headtohead} gives Breast Cancer (0.975), Adult (0.858),
Covertype (0.733) - the same order.  Where classes are recoverable
without labels, they are easy to recover with them.

% - Discussion -
\section{Discussion}
\label{sec:discussion}

\subsection{Answering the central question}
Our evidence says the honest answer is narrower than the question implies.

\textbf{Boosting wins when the base learner's bias is the binding
constraint and the labels are trustworthy.}  Breast Cancer is that case:
clean, balanced, low-dimensional enough that a stump is informative, and
AdaBoost takes it (0.9754 vs 0.9595).  Table~\ref{tab:biasvar} says why -
AdaBoost is the only model here that reduced bias.  When variance is
already handled and what remains is systematic error, only boosting has a
mechanism for it.

\textbf{Bagging wins when variance dominates, when labels are noisy, or
when the weak learner is too weak.}  All three appear in our results: RF
loses 0.9 points to AdaBoost's 6.2 under 20\% noise, and on Covertype,
where a stump cannot address seven classes, AdaBoost (0.5885) finishes
below a single tree while RF wins (0.7326).

\textbf{Most of the time it does not matter much.}  On Adult the two differ
by 0.0005 against a standard deviation of 0.01.  This is worth stating
plainly, because a study of this shape invites you to declare a winner from
a difference the folds cannot resolve.  The defensible claim is not that
one method is better but that they fail differently, and predictably from
Table~\ref{tab:biasvar}.

If forced to pick a default: Random Forest - not because it won more
often, since it did not, but because its failure mode is graceful.  It
cannot overfit as trees are added, its bad case is mediocrity rather than
collapse, and OOB gives a free honest estimate of how it is doing.
AdaBoost is the better instrument when the labels are clean and bias is
what hurts.

\subsection{Dataset characteristics}
\textbf{Dimensionality} hurt the unsupervised half far more than the
supervised half.  Adult's 104 one-hot features left DBSCAN at chance and
needed 74 PCs for 90\% variance, yet both ensembles reached about 0.857.
Trees do implicit feature selection - they never split on a useless
indicator - while distance-based methods must use every dimension given.

\textbf{Class imbalance} was the most instructive.  Covertype's 0.47\%
minority is invisible in accuracy and obvious in macro-F1, where RF drops
from 0.7326 accuracy to 0.4632.  Reporting without resampling is the right
default for a comparative study, but it means our Covertype numbers
describe models that largely ignore the rare classes.  Class weights map
naturally onto the \texttt{sample\_weight} interface our tree already
exposes, which makes them a better fit here than SMOTE - interpolating
synthetic neighbours between one-hot indicators produces points
corresponding to no real category assignment.  Quantifying that trade-off
is the most obvious thing missing from this study.

\textbf{Noise} produced the clearest separation, for no mysterious reason:
equation~(\ref{eq:update}) cannot tell a hard example from a wrong one.

\subsection{Do clusters predict classifier performance?}
Yes, on our three datasets, and the ordering is exact.  ARI measures
whether class structure lies along the directions of greatest variance,
which is what unsupervised methods can see; when it does, the problem is
easy for everything.  When it does not - Adult, where classes depend on
interactions between sparse indicators rather than dominant variance
directions - clustering finds nothing while the ensembles still reach
0.857, because a tree carves out structure no centroid distance reveals.
So the relationship is one-directional: high ARI predicts easy
classification, but low ARI does not predict hard classification.  Three
datasets is also three points, which cannot support more than a suggestive
trend.

\subsection{Limitations}
Subsampling is the main one: capping Adult at 8{,}000 rows and the sweeps
at 2{,}000 depresses every absolute number, and the caps were set by our
implementation's speed rather than by the science.  Model rankings are
unaffected, since all models saw the same rows, but these are not the best
achievable accuracies here.  The bias-variance decomposition runs on one
dataset, as the brief specifies, so its clean confirmation should not be
over-generalised.  Our AdaBoost is discrete SAMME; SAMME.R would likely
improve the Covertype probabilities given the AUC/accuracy gap.  And the
Covertype noise anomaly deserves proper investigation rather than the
plausible story we offered.

% - Conclusion -
\section{Conclusion}
\label{sec:conclusion}
We implemented a CART decision tree, AdaBoost with stumps, and a Random
Forest from scratch, verified them against sklearn (within 0.6\% on the two
larger datasets), and used them to test when boosting beats bagging.

The bias-variance decomposition gave the sharpest answer: bagging left the
base learner's bias exactly unchanged while cutting variance by 64\%, and
AdaBoost was the only method to reduce bias.  Everything else follows.
Boosting won on clean, balanced Breast Cancer, where bias was what remained
to remove.  Bagging won decisively under label noise (0.9 points lost
against AdaBoost's 6.2 at $\eta=0.20$) and on multi-class Covertype, where
a stump is too weak to boost and AdaBoost finished below a single tree.  On
Adult the two were indistinguishable.  The unsupervised pipeline ranked the
datasets by ARI in the same order the classifiers ranked them by accuracy,
suggesting the two views agree about which problems are hard.

The most useful extension is the imbalance treatment we implemented but did
not evaluate: measuring the accuracy/macro-F1 trade-off from class weights
on Covertype would address this study's weakest results directly.  Beyond
that, Gradient Boosting would separate AdaBoost's exponential loss from
boosting in general, and t-SNE would test whether Adult's class structure
is genuinely absent from low dimensions or merely invisible to a linear
projection.

% - Tools & Acknowledgements -
\section*{Tools and Acknowledgements}
In line with the course's academic integrity policy, we disclose that an AI
assistant Claude Sonnet 4.5 was used during this project.  It helped to write:
test cases, format plots and report; 'fit', 'compute oob score' in random forest; 'fit' in adaboost, 'best split' and 'build' in decision tree, dbscan, kmeans 

% ===== BIBLIOGRAPHY ===================================================
\begin{thebibliography}{00}

\bibitem{hastie2009elements}
T.~Hastie, R.~Tibshirani, and J.~Friedman,
\textit{The Elements of Statistical Learning}, 2nd ed.
Springer, 2009.
[Online]. Available: \url{https://hastie.su.domains/ElemStatLearn/}

\bibitem{quinlan1986induction}
J.~R.~Quinlan,
``Induction of decision trees,''
\textit{Machine Learning}, vol.~1, no.~1, pp.~81-106, 1986.

\bibitem{freund1997decision}
Y.~Freund and R.~E.~Schapire,
``A decision-theoretic generalization of on-line learning and an
application to boosting,''
\textit{J. Comput. Syst. Sci.}, vol.~55, no.~1, pp.~119-139, 1997.

\bibitem{breiman1996bagging}
L.~Breiman,
``Bagging predictors,''
\textit{Machine Learning}, vol.~24, no.~2, pp.~123-140, 1996.

\bibitem{breiman2001random}
L.~Breiman,
``Random forests,''
\textit{Machine Learning}, vol.~45, no.~1, pp.~5-32, 2001.

\bibitem{pearson1901lines}
K.~Pearson,
``On lines and planes of closest fit to systems of points in space,''
\textit{Philosophical Magazine}, vol.~2, no.~11, pp.~559-572, 1901.

\bibitem{lloyd1982least}
S.~P.~Lloyd,
``Least squares quantization in PCM,''
\textit{IEEE Trans. Inf. Theory}, vol.~28, no.~2, pp.~129-137, 1982.

\bibitem{ester1996density}
M.~Ester, H.-P.~Kriegel, J.~Sander, and X.~Xu,
``A density-based algorithm for discovering clusters in large
spatial databases with noise,''
in \textit{Proc. 2nd ACM SIGKDD}, 1996, pp.~226-231.

\end{thebibliography}

% ===== APPENDICES (optional) ==========================================
% \appendices
% \section{Code snippets}
% \lstinputlisting[style=pythonstyle]{../src/trees/decision_tree.py}

\end{document}
